\chapter{2005年全国II卷（理）}
\section{单选题}
\begin{problem}
函数 \( f(x) = |\sin x + \cos x| \) 的最小正周期是（\quad）

\choices
    {\( \frac{\pi}{4} \)}
    {\( \frac{\pi}{2} \)}
    {\(\pi\)}
    {\(2\pi\)}
\end{problem}

\begin{answer}
C
\end{answer}

\begin{solution}
因为 \(f(x)=\sqrt{2}\left|\sin\left(x+\frac{\pi}{4}\right)\right|\)，所以 \(f(x+\pi)=f(x)\)，即 \(\pi\) 是它的一个正周期. 函数的零点为 \(x=-\frac{\pi}{4}+k\pi\)（\(k\in\Z\)）. 若正数 \(T\) 是它的周期，则 \(-\frac{\pi}{4}+T\) 也必须是零点，从而 \(T=k\pi\)（\(k\in\Z\)）. 因为 \(T>0\)，所以 \(T\geqslant\pi\)，故最小正周期为 \(\pi\)，选 C.
\end{solution}

\begin{problem}
正方体 \(ABCD - A_{1}B_{1}C_{1}D_{1}\) 中，\(P\)，\(Q\)，\(R\) 分别是 \(AB\)，\(AD\)，\(B_{1}C_{1}\) 的中点. 那么，正方体的过 \(P\)，\(Q\)，\(R\) 的截面图形是（\quad）

\choices
    {三角形}
    {四边形}
    {五边形}
    {六边形}
\end{problem}

\begin{answer}
D
\end{answer}

\begin{solution}
设正方体棱长为 \(1\)，取 \(A(0,0,0)\)，\(B(1,0,0)\)，\(D(0,1,0)\)，\(A_1(0,0,1)\). 则 \(P\left(\frac{1}{2},0,0\right)\)，\(Q\left(0,\frac{1}{2},0\right)\)，\(R\left(1,\frac{1}{2},1\right)\). 由 \(\overrightarrow{PQ}=\left(-\frac{1}{2},\frac{1}{2},0\right)\)，\(\overrightarrow{PR}=\left(\frac{1}{2},\frac{1}{2},1\right)\)，可取截面平面的法向量为 \((1,1,-1)\)，故截面平面方程为 \(x+y-z=\frac{1}{2}\). 该平面与正方体的六条棱 \(AB\)、\(AD\)、\(BB_1\)、\(B_1C_1\)、\(C_1D_1\)、\(DD_1\) 分别相交，交点中包括已知的 \(P\)、\(Q\)、\(R\)，其余三个交点分别位于 \(BB_1\)、\(C_1D_1\)、\(DD_1\) 上. 因为截面平面没有经过正方体顶点，六个交点依次构成截面的六个顶点，所以截面为六边形，选 D.
\end{solution}

\begin{problem}
函数 \( y = \sqrt[3]{x^2 - 1} \)（\(x \leqslant 0\)）的反函数是（\quad）

\choices
    {\( y = \sqrt{(x + 1)^3} \)（\(x \geqslant -1\)）}
    {\( y = -\sqrt{(x + 1)^3} \)（\(x \geqslant -1\)）}
    {\( y = \sqrt{(x + 1)^3} \)（\(x \geqslant 0\)）}
    {\( y = -\sqrt{(x + 1)^3} \)（\(x \geqslant 0\)）}
\end{problem}

\begin{answer}
B
\end{answer}

\begin{solution}
在 \(x\leqslant0\) 上，函数 \(x^2\) 随 \(x\) 增大而减小，立方根函数严格递增，所以原函数是一一对应的. 由 \(y=\sqrt[3]{x^2-1}\) 得 \(x^2=y^3+1\). 原函数的定义域限制为 \(x\leqslant0\)，所以反解时应取负根，即 \(x=-\sqrt{y^3+1}\). 原函数的值域为 \(y\geqslant-1\)，交换 \(x\)，\(y\) 后得到反函数 \(y=-\sqrt{(x+1)^3}\)，其定义域为 \(x\geqslant-1\)，选 B.
\end{solution}

\begin{problem}
已知函数 \( y = \tan \omega x \) 在 \( \left(-\frac{\pi}{2}, \frac{\pi}{2}\right) \) 内是减函数，则（\quad）

\choices
    {\(0 < \omega \leqslant 1\)}
    {\(-1 \leqslant \omega < 0\)}
    {\(\omega \geqslant 1\)}
    {\(\omega \leqslant -1\)}
\end{problem}

\begin{answer}
B
\end{answer}

\begin{solution}
函数在整个 \(\left(-\frac{\pi}{2},\frac{\pi}{2}\right)\) 内有定义. 若 \(|\omega|>1\)，取 \(x_0=\frac{\pi}{2\omega}\)，则 \(x_0\in\left(-\frac{\pi}{2},\frac{\pi}{2}\right)\) 且 \(\omega x_0=\frac{\pi}{2}\)，函数无定义，所以必须有 \(|\omega|\leqslant1\). 在此条件下，\(\tan(\omega x)\) 在该区间内有定义，且 \(\frac{\mathrm{d}}{\mathrm{d}x}\tan(\omega x)=\omega\sec^2(\omega x)\)，其中 \(\sec^2(\omega x)>0\). 因而严格减函数要求 \(\omega<0\)，合并得 \(-1\leqslant\omega<0\)，选 B.
\end{solution}

\begin{problem}
设 \(a\)，\(b\)，\(c\)，\(d \in \R\)，若 \(\frac{a + b\i}{c + d\i}\) 为实数，则（\quad）

\choices
    {\(b c + a d \neq 0\)}
    {\(bc - ad \neq 0\)}
    {\(bc - ad = 0\)}
    {\(bc + ad = 0\)}
\end{problem}

\begin{answer}
C
\end{answer}

\begin{solution}
因为分式有意义，故 \(c^2+d^2\neq0\). 分母实数化得
\[
\frac{a+b\i}{c+d\i}=\frac{(a+b\i)(c-d\i)}{c^2+d^2}=\frac{ac+bd+(bc-ad)\i}{c^2+d^2}
\]
分式为实数当且仅当虚部为零，即 \(bc-ad=0\)，选 C.
\end{solution}

\begin{problem}
已知双曲线 \(\frac{x^2}{6} - \frac{y^2}{3} = 1\) 的焦点为 \(F_1\)，\(F_2\)，点 \(M\) 在双曲线上且 \(MF_1 \perp x\) 轴，则 \(F_1\) 到直线 \(F_2M\) 的距离为（\quad）

\choices
    {\(\frac{3\sqrt{6}}{5}\)}
    {\(\frac{5\sqrt{6}}{6}\)}
    {\(\frac{6}{5}\)}
    {\(\frac{5}{6}\)}
\end{problem}

\begin{answer}
C
\end{answer}

\begin{solution}
双曲线的焦半距满足 \(c^2=6+3=9\)，故可取 \(F_1(-3,0)\)，\(F_2(3,0)\). 由 \(MF_1\perp x\) 轴知 \(M\) 的横坐标为 \(-3\). 代入双曲线方程得 \(\frac{9}{6}-\frac{y_M^2}{3}=1\)，所以 \(y_M^2=\frac{3}{2}\).
于是 \(\lvert F_1M\rvert=\frac{\sqrt{6}}{2}\)，\(\lvert F_2M\rvert=\sqrt{6^2+\left(\frac{\sqrt{6}}{2}\right)^2}=\frac{5\sqrt{6}}{2}\). 三角形 \(F_1F_2M\) 的面积为 \(\frac{1}{2}\cdot6\cdot\frac{\sqrt{6}}{2}=\frac{3\sqrt{6}}{2}\). 设所求距离为 \(h\)，则 \(\frac{1}{2}\lvert F_2M\rvert h=\frac{3\sqrt{6}}{2}\)，解得 \(h=\frac{6}{5}\)，选 C.
\end{solution}

\begin{problem}
锐角三角形的内角 \(A\)，\(B\) 满足 \(\tan A - \frac{1}{\sin 2A} = \tan B\)，则有（\quad）

\choices
    {\(\sin 2A - \cos B = 0\)}
    {\(\sin 2A + \cos B = 0\)}
    {\(\sin 2A - \sin B = 0\)}
    {\(\sin 2A + \sin B = 0\)}
\end{problem}

\begin{answer}
A
\end{answer}

\begin{solution}
由 \(\tan A-\frac{1}{\sin2A}=\frac{2\sin^2A-1}{2\sin A\cos A}=-\cot2A\)，得 \(\tan B=-\cot2A\). 因为三角形为锐角三角形且 \(\tan B>0\)，所以 \(\frac{\pi}{2}<2A<\pi\)，从而 \(B=2A-\frac{\pi}{2}\). 因此 \(\cos B=\cos\left(2A-\frac{\pi}{2}\right)=\sin2A\)，即 \(\sin2A-\cos B=0\)，选 A.
\end{solution}

\begin{problem}
已知点 \(A(\sqrt{3}, 1)\)，\(B(0, 0)\)，\(C(\sqrt{3}, 0)\). 设 \(\angle BAC\) 的平分线 \(AE\) 与 \(BC\) 相交于 \(E\)，那么有 \(\overrightarrow{BC} = \lambda \overrightarrow{CE}\)，其中 \(\lambda\) 等于（\quad）

\choices
    {\(2\)}
    {\(\frac{1}{2}\)}
    {\(-3\)}
    {\(-\frac{1}{3}\)}
\end{problem}

\begin{answer}
C
\end{answer}

\begin{solution}
\(AB=\sqrt{(\sqrt{3})^2+1^2}=2\)，而 \(AC=1\). 由角平分线定理，\(\frac{BE}{CE}=\frac{AB}{AC}=2\)，又 \(BE+CE=BC=\sqrt{3}\)，所以 \(CE=\frac{\sqrt{3}}{3}\). 方向上 \(\overrightarrow{BC}\) 与 \(\overrightarrow{CE}\) 相反，且 \(\lvert BC\rvert=3\lvert CE\rvert\)，故 \(\overrightarrow{BC}=-3\overrightarrow{CE}\)，即 \(\lambda=-3\)，选 C.
\end{solution}

\begin{problem}
已知集合 \(M = \{x \mid x^2 - 3x - 28 \leqslant 0\}\)，\(N = \{x \mid x^2 - x - 6 > 0\}\)，则 \(M \cap N\) 为（\quad）

\choices
    {\(\left\{x\mid -4\leqslant x < -2\text{\ 或\ }3 < x\leqslant 7\right\}\)}
    {\(\left\{x\mid -4 < x\leqslant -2\text{\ 或\ }3\leqslant x < 7\right\}\)}
    {\(\left\{x\mid x\leqslant -2\text{\ 或\ }x > 3\right\}\)}
   {\(\left\{x\mid x < -2\text{\ 或\ }x\geqslant 3\right\}\)}
\end{problem}

\begin{answer}
A
\end{answer}

\begin{solution}
\(x^2-3x-28=(x-7)(x+4)\)，所以 \(M=[-4,7]\). 又 \(x^2-x-6=(x-3)(x+2)>0\)，所以 \(N=(-\infty,-2)\cup(3,+\infty)\). 因而 \(M\cap N=[-4,-2)\cup(3,7]\)，选 A.
\end{solution}

\begin{problem}
点 \(P\) 在平面上作匀速直线运动，速度向量 \(\bs{v}=(4,-3)\)（即点 \(P\) 的运动方向与 \(\bs{v}\) 相同，且每秒移动的距离为 \(|\bs{v}|\) 个单位）. 设开始时点 \(P\) 的坐标为 \((-10,10)\)，则 \(5\) 秒后点 \(P\) 的坐标为（\quad）

\choices
    {\((-2, 4)\)}
    {\((-30, 25)\)}
    {\((10, -5)\)}
    {\((5, -10)\)}
\end{problem}

\begin{answer}
C
\end{answer}

\begin{solution}
\(\lvert\bs v\rvert=\sqrt{4^2+(-3)^2}=5\)，所以每秒位移向量就是 \(\bs v\). 五秒后的位移为 \(5\bs v=(20,-15)\)，故点 \(P\) 的坐标为 \((-10,10)+(20,-15)=(10,-5)\)，选 C.
\end{solution}

\begin{problem}
如果 \(a_{1}\)，\(a_{2}\)，\(\cdots\)，\(a_{8}\) 为各项都大于零的等差数列，公差 \(d \neq 0\)，则（\quad）

\choices
    {\(a_{1}a_{8} > a_{4}a_{5}\)}
    {\(a_{1}a_{8} < a_{4}a_{5}\)}
    {\(a_{1} + a_{8} > a_{4} + a_{5}\)}
    {\(a_{1}a_{8} = a_{4}a_{5}\)}
\end{problem}

\begin{answer}
B
\end{answer}

\begin{solution}
设公差为 \(d\)，则 \(a_4=a_1+3d\)，\(a_5=a_1+4d\)，\(a_8=a_1+7d\). 于是
\[
a_4a_5-a_1a_8=(a_1+3d)(a_1+4d)-a_1(a_1+7d)=12d^2>0
\]
因为 \(d\neq0\). 所以 \(a_1a_8<a_4a_5\)，选 B；同时 \(a_1+a_8=a_4+a_5\)，故其他选项均不成立.
\end{solution}

\begin{problem}
将半径都为 \(1\) 的 \(4\) 个铅球完全装入形状为正四面体的容器里，这个正四面体的高最小值为（\quad）

\choices
    {\(\frac{\sqrt{3} + 2\sqrt{6}}{3}\)}
    {\(2 + \frac{2\sqrt{6}}{3}\)}
    {\(4 + \frac{2\sqrt{6}}{3}\)}
    {\(\frac{4\sqrt{3} + 2\sqrt{6}}{3}\)}
\end{problem}

\begin{answer}
C
\end{answer}

\begin{solution}
设外部正四面体的高为 \(H\). 其内切球半径为 \(H/4\). 每个半径为 \(1\) 的球完全在容器内，等价于它的球心到外部正四面体的每个面的距离都不小于 \(1\). 因此，四个球心都位于将外部正四面体的四个面分别向内平移 \(1\) 所得到的正四面体 \(T\) 中. 由于内切半径减少了 \(1\)，\(T\) 的高为
\(h=H-4\).
设 \(T\) 的棱长为 \(l\). 正四面体的高与棱长满足 \(h=l\sqrt{\frac{2}{3}}\)，而 \(T\) 的直径就是棱长 \(l\). 四个球互不穿透，所以任意两个球心的距离都不小于 \(2\)，从而 \(l\geqslant2\). 因此
\[
H-4=h=l\sqrt{\frac{2}{3}}\geqslant2\sqrt{\frac{2}{3}}=\frac{2\sqrt{6}}{3}
\]
当 \(T\) 的四个顶点作为四个球心时，四个球心两两距离均为 \(2\)，且每个球到外部正四面体的相应三个面均相切，于是下界可以达到. 所以外部正四面体的最小高为
\(H=4+\frac{2\sqrt{6}}{3}\).
选 C.
\end{solution}

\section{填空题}
\begin{problem}
圆心为 \((1,2)\) 且与直线 \( 5x - 12y - 7 = 0 \) 相切的圆的方程为\fillinblank{}.
\end{problem}

\begin{answer}
\((x-1)^2+(y-2)^2=4\)
\end{answer}

\begin{solution}
圆心到切线的距离就是半径，故
\[
r=\frac{|5\cdot1-12\cdot2-7|}{\sqrt{5^2+(-12)^2}}=\frac{26}{13}=2
\]
所以圆的方程为 \((x-1)^2+(y-2)^2=4\).
\end{solution}

\begin{problem}
设 \(\alpha\) 为第四象限的角，若 \(\frac{\sin 3\alpha}{\sin \alpha} = \frac{13}{5}\)，则 \(\tan 2\alpha =\)\fillinblank{}.
\end{problem}

\begin{answer}
\(-\frac{3}{4}\)
\end{answer}

\begin{solution}
由三倍角公式，\(\frac{\sin3\alpha}{\sin\alpha}=3-4\sin^2\alpha=\frac{13}{5}\)，所以 \(\sin^2\alpha=\frac{1}{10}\)，且因 \(\alpha\) 在第四象限，\(\cos^2\alpha=\frac{9}{10}\)，\(\tan\alpha=-\frac{1}{3}\). 因而
\[
\tan2\alpha=\frac{2\tan\alpha}{1-\tan^2\alpha}=\frac{-2/3}{1-1/9}=-\frac{3}{4}
\]
\end{solution}

\begin{problem}
在由数字 \(0\)，\(1\)，\(2\)，\(3\)，\(4\)，\(5\) 所组成的没有重复数字的四位数中，不能被 \(5\) 整除的数共有\fillinblank{} 个.
\end{problem}

\begin{answer}
192
\end{answer}

\begin{solution}
四位数的首位不能为 \(0\)，总数为 \(5\cdot\mathrm A_{5}^{3}=300\). 能被 \(5\) 整除的数末位只能是 \(0\) 或 \(5\). 末位为 \(0\) 时有 \(5\cdot\mathrm A_{4}^{2}=60\) 个；末位为 \(5\) 时首位有 \(1\)，\(2\)，\(3\)，\(4\) 四种选择，另两位有 \(\mathrm A_{4}^{2}=12\) 种排法，共 \(4\cdot12=48\) 个. 所以不能被 \(5\) 整除的数有 \(300-60-48=192\) 个.
\end{solution}

\begin{problem}
下面是关于三棱锥的四个命题：

\begin{circlelist}
\item 底面是等边三角形，侧面与底面所成的二面角都相等的三棱锥是正三棱锥；
\item 底面是等边三角形，侧面都是等腰三角形的三棱锥是正三棱锥；
\item 底面是等边三角形，侧面的面积都相等的三棱锥是正三棱锥；
\item 侧棱与底面所成的角都相等，且侧面与底面所成的二面角都相等的三棱锥是正三棱锥.
\end{circlelist}

其中，真命题的编号是\fillinblank{}. （写出所有真命题的编号）
\end{problem}

\begin{answer}
\circled{1}，\circled{4}
\end{answer}

\begin{solution}
设顶点为 \(P\)，\(H\) 为 \(P\) 在底面上的射影.

对于命题 1，在垂直于每条底边的截面内考察内部二面角. 设 \(PH=h>0\)，并把 \(H\) 到某条底边所在直线、朝向底面内部一侧的有向距离记为 \(\delta\). 截面中侧面内射线的方向可表示为 \((\delta,h)\)，所以该二面角 \(\varphi\) 满足 \(\cot\varphi=\delta/h\). 因为 \(h\) 固定，三个二面角相等就等价于 \(H\) 到三条底边所在直线的有向距离相等. 对等边三角形，三个有向距离之和恒为三倍内切圆半径，故它们相等时 \(H\) 恰为内心，也就是中心，从而命题 1 正确.

命题 2 不正确. 例如取底面为单位等边三角形 \(A(0,0,0)\)，\(B(1,0,0)\)，\(C\left(\frac12,\frac{\sqrt3}{2},0\right)\)，取 \(P\left(\frac12,0,\frac{\sqrt3}{2}\right)\). 此时 \(PA=PB=AB=BC=CA=1\)，所以三个侧面分别因 \(PA=PB\)、\(PB=BC\)、\(PA=CA\) 而都是等腰三角形，但 \(PC=\sqrt{\frac32}\neq1\)，故不是正三棱锥.

命题 3 不正确. 取 \(H\) 为底面等边三角形的一个旁心，取 \(P\) 在过 \(H\) 的底面法线上且不在底面内. 设 \(H\) 到三条底边所在直线的距离的绝对值均为 \(r\)，且 \(PH=h\). 则从 \(P\) 到三条底边所在直线的距离均为 \(\sqrt{h^2+r^2}\)，所以三个侧面的高相等，三个侧面的面积也相等；但 \(H\) 不是底面中心，故三棱锥不是正三棱锥.

对于命题 4，设侧棱与底面所成的角为 \(\theta\)，则 \(\sin\theta=PH/PA=PH/PB=PH/PC\)，从而 \(PA=PB=PC\). 因为 \(A\)，\(B\)，\(C\) 是等边三角形的顶点，平面内到三点等距的点只有其中心，所以 \(H\) 是底面中心，三棱锥为正三棱锥. 因此真命题的编号为 1，4.
\end{solution}

\section{解答题}
\begin{problem}
设函数 \( f(x) = 2^{|x + 1| - |x - 1|} \)，求使 \( f(x) \geqslant 2\sqrt{2} \) 的 \( x \) 的取值范围.
\end{problem}

\begin{solution}
因为底数 \(2>1\)，不等式等价于 \(|x+1|-|x-1|\geqslant\log_2(2\sqrt2)=\frac32\).
分三段讨论：当 \(x\leqslant-1\) 时，左端为 \(-2\)；当 \(-1\leqslant x\leqslant1\) 时，左端为 \(2x\)；当 \(x\geqslant1\) 时，左端为 \(2\). 因而中间一段给出 \(x\geqslant\frac34\)，右侧一段全部满足，左侧一段无解. 所以 \(x\in\left[\frac34,+\infty\right)\).
\end{solution}

\begin{problem}
已知 \( \{a_{n}\} \) 是各项均为正数的等差数列，\(\lg a_1\)，\(\lg a_2\)，\(\lg a_4\) 成等差数列. 又 \(b_{n} = \frac{1}{a_{2^n}}\)，\(n = 1\)，\(2\)，\(3\)，\(\cdots\).

\begin{enumerate}
\item 证明 \( \{b_{n}\} \) 为等比数列；
\item 如果无穷等比数列 \( \{b_{n}\} \) 各项的和 \( S = \frac{1}{3} \)，求数列 \( \{a_{n}\} \) 的首项 \( a_{1} \) 和公差 \( d \).
\end{enumerate}

注：无穷数列各项的和即当 \(n \to \infty\) 时数列前 \(n\) 项和的极限.
\end{problem}

\begin{solution}
设等差数列的公差为 \(d\). 由 \(\lg a_1\)，\(\lg a_2\)，\(\lg a_4\) 成等差数列，得 \(2\lg a_2=\lg a_1+\lg a_4\)，从而 \(a_2^2=a_1a_4\). 代入 \(a_2=a_1+d\)，\(a_4=a_1+3d\)，得到 \((a_1+d)^2=a_1(a_1+3d)\)，即 \(d(d-a_1)=0\).

若 \(d=0\)，则 \(a_{2^n}=a_1\)，所以 \(b_n=\frac1{a_1}\)，这是公比为 \(1\) 的等比数列. 若 \(d=a_1\)，则 \(a_n=na_1\)，从而 \(a_{2^n}=2^n a_1\)，\(b_n=\frac{1}{2^n a_1}=\frac{1}{2a_1}\left(\frac12\right)^{n-1}\)，这是公比为 \(\frac12\) 的等比数列. 因此 \(\{b_n\}\) 必为等比数列.

第二问中，无穷等比数列的和有限，排除公比为 \(1\) 的情形，故 \(d=a_1\). 此时 \(b_1=\frac{1}{2d}\)，公比为 \(\frac12\)，所以
\[
S=\frac{1/(2d)}{1-1/2}=\frac1d=\frac13
\]
故 \(d=3\)，且 \(a_1=d=3\).
\end{solution}

\begin{problem}
甲，乙两队进行一场排球比赛. 根据以往经验，单局比赛甲队胜乙队的概率为 \(0.6\). 本场比赛采用五局三胜制，即先胜三局的队获胜，比赛结束. 各局比赛相互间没有影响. 令 \(\xi\) 为本场比赛的局数. 求 \(\xi\) 的概率分布和数学期望. （精确到 \(0.0001\)）
\end{problem}

\begin{solution}
设甲队单局获胜概率为 \(p=0.6\)，乙队单局获胜概率为 \(q=0.4\). 比赛恰好进行三局，当且仅当某队连续胜三局，因此
\[
P(\xi=3)=p^3+q^3=0.216+0.064=0.2800
\]
比赛恰好进行四局时，前四局中获胜队在前四局恰好取得第三场胜利，故
\[
P(\xi=4)=\mathrm C_3^2p^3q+\mathrm C_3^2q^3p=3(0.6^3\cdot0.4+0.4^3\cdot0.6)=0.3744
\]
比赛恰好进行五局，当且仅当前四局双方各胜两局，所以 \(P(\xi=5)=\mathrm C_4^2p^2q^2=6\cdot0.6^2\cdot0.4^2=0.3456\).
三种概率之和为 \(1\)，故 \(\xi\) 的概率分布为上述三项，其他取值的概率为 \(0\). 数学期望为
\[
E\xi=3\cdot0.2800+4\cdot0.3744+5\cdot0.3456=4.0656
\]
\end{solution}

\begin{problem}
如图，四棱锥 \(P - ABCD\) 中，底面 \(ABCD\) 为矩形，\(PD \perp\) 底面 \(ABCD\)，\(AD = PD\)，\(E\)，\(F\) 分别为 \(CD\)，\(PB\) 的中点.

\begin{enumerate}
\item 求证：\(EF \perp\) 平面 \(PAB\)；
\item 设 \(AB = \sqrt{2}BC\)，求 \(AC\) 与平面 \(AEF\) 所成的角的大小.
\end{enumerate}

\bitmapfigure[max width=0.45\linewidth]{img/2005/national_paper_2_science/q20_fig1.png}
\end{problem}

\begin{solution}
建立空间直角坐标系，令 \(D(0,0,0)\)，\(A(a,0,0)\)，\(B(a,b,0)\)，\(C(0,b,0)\)，其中 \(a=AD=PD\)，且 \(P(0,0,a)\). 则 \(E\left(0,\frac b2,0\right)\)，\(F\left(\frac a2,\frac b2,\frac a2\right)\).
于是 \(\overrightarrow{EF}=\left(\frac a2,0,\frac a2\right)\)，\(\overrightarrow{AB}=(0,b,0)\)，\(\overrightarrow{AP}=(-a,0,a)\). 有 \(\overrightarrow{EF}\cdot\overrightarrow{AB}=0\)，\(\overrightarrow{EF}\cdot\overrightarrow{AP}=0\)，而 \(AB\) 与 \(AP\) 是平面 \(PAB\) 内相交的两条直线，所以 \(EF\perp\) 平面 \(PAB\).

再由 \(AB=\sqrt2BC\) 得 \(b=\sqrt2a\). 取平面 \(AEF\) 内的两个向量 \(\overrightarrow{AE}=\left(-a,\frac b2,0\right)\) 和 \(\overrightarrow{AF}=\left(-\frac a2,\frac b2,\frac a2\right)\).
可取该平面的法向量为 \(\bs{n}=(b,2a,-b)\)，而 \(\overrightarrow{AC}=(-a,b,0)\). 设直线 \(AC\) 与平面 \(AEF\) 所成的角为 \(\theta\)，则
\[
\sin\theta=\frac{|\overrightarrow{AC}\cdot\bs{n}|}{|\overrightarrow{AC}|\,|\bs{n}|}=\frac{ab}{\sqrt{a^2+b^2}\sqrt{2b^2+4a^2}}=\frac{\sqrt3}{6}
\]
因此 \(\theta=\arcsin\frac{\sqrt3}{6}\).
\end{solution}

\begin{problem}
\(P\)，\(Q\)，\(M\)，\(N\) 四点都在椭圆 \(x^{2} + \frac{y^{2}}{2} = 1\) 上，\(F\) 为椭圆在 \(y\) 轴正半轴上的焦点. 已知 \(\overrightarrow{PF}\) 与 \(\overrightarrow{FQ}\) 共线，\(\overrightarrow{MF}\) 与 \(\overrightarrow{FN}\) 共线，且 \(\overrightarrow{PF} \cdot \overrightarrow{MF} = 0\). 求四边形 \(PMQN\) 的面积的最小值和最大值.
\end{problem}

\begin{solution}
椭圆可写为 \(\frac{x^2}{1^2}+\frac{y^2}{(\sqrt2)^2}=1\)，故焦点为 \(F(0,1)\). 由题意，弦 \(PQ\) 与弦 \(MN\) 都过 \(F\)，且互相垂直，因此四边形 \(PMQN\) 的两条对角线为 \(PQ\)、\(MN\)，其面积为 \(S=\frac12|PQ||MN|\).

先设直线 \(PQ\) 的斜率为非零有限值 \(k\). 直线方程为 \(y=1+kx\)，代入椭圆方程得
\[
(k^2+2)x^2+2kx-1=0
\]
该方程的判别式为 \(\Delta=4k^2+4(k^2+2)=8(k^2+1)\)，所以两交点横坐标之差为 \(\frac{\sqrt{\Delta}}{k^2+2}=\frac{2\sqrt{2(k^2+1)}}{k^2+2}\)，进而
\[
|PQ|=\frac{2\sqrt2(k^2+1)}{k^2+2}
\]
垂线 \(MN\) 的斜率为 \(-1/k\). 将上面得到的弦长公式中的斜率替换为 \(-1/k\)，得到
\[
|MN|=\frac{2\sqrt2(k^2+1)}{2k^2+1}
\]
令 \(u=k^2\geqslant0\)，则
\[
S=\frac{4(1+u)^2}{(u+2)(2u+1)}
\]
由
\[
2(u+2)(2u+1)-4(1+u)^2=2u\geqslant0
\]
得 \(S\leqslant2\)；又
\[
36(1+u)^2-16(u+2)(2u+1)=4(u-1)^2\geqslant0
\]
得 \(S\geqslant\frac{16}{9}\). 当 \(u=1\) 时取到下界. 对遗漏的退化斜率，若 \(k=0\)，则 \(PQ\) 为直线 \(y=1\)，\(MN\) 为直线 \(x=0\)，故 \(|PQ|=\sqrt{2}\)，\(|MN|=2\sqrt{2}\)，仍有 \(S=2\)；若 \(PQ\) 为竖直弦，则 \(MN\) 为水平弦，同样得到 \(S=2\). 因此最大值也能取到，面积最小值为 \(\frac{16}{9}\)，最大值为 \(2\).
\end{solution}

\begin{problem}
已知 \(a \geqslant 0\)，函数 \(f(x) = (x^2 - 2ax)\e^x\).

\begin{enumerate}
\item 当 \(x\) 为何值时，\(f(x)\) 取得最小值？ 证明你的结论；
\item 设 \(f(x)\) 在 \([-1,1]\) 上是单调函数，求 \(a\) 的取值范围.
\end{enumerate}
\end{problem}

\begin{solution}
\begin{enumerate}
\item
\(f'(x)=\bigl[x^2+2(1-a)x-2a\bigr]\e^x\). 令 \(r_1=a-1-\sqrt{a^2+1}\)，\(r_2=a-1+\sqrt{a^2+1}\)，则 \(r_1<r_2\)，且 \(x^2+2(1-a)x-2a=(x-r_1)(x-r_2)\).
因为 \(\e^x>0\)，所以 \(f\) 在 \((-\infty,r_1)\) 上递增，在 \((r_1,r_2)\) 上递减，在 \((r_2,+\infty)\) 上递增. 又当 \(x\to-\infty\) 时 \(f(x)\to0\)，当 \(x\to+\infty\) 时 \(f(x)\to+\infty\). 当 \(a>0\) 时，\(0<r_2<2a\)，故 \(f(r_2)<0\)；当 \(a=0\) 时，\(r_2=0\)，且 \(f(x)=x^2\e^x\geqslant0=f(0)\). 所以无论哪种情形，函数都在 \(x=a-1+\sqrt{a^2+1}\) 时取得全局最小值.
\item
在 \([-1,1]\) 上，\(f'(x)\) 的符号由 \(g(x)=x^2+2(1-a)x-2a\) 决定. 因为 \(g(-1)=-1<0\)，不可能在整个区间上满足 \(g(x)\geqslant0\)，所以若 \(f\) 单调，必须有 \(g(x)\leqslant0\) 对一切 \(x\in[-1,1]\) 成立. 二次函数 \(g\) 开口向上，在闭区间上的最大值为端点值，而 \(g(-1)=-1\)，\(g(1)=3-4a\).
因此充要条件为 \(3-4a\leqslant0\)，即 \(a\geqslant\frac34\).
\end{enumerate}
\end{solution}
